\begin{pastexam}[label={px:ntut-113-ci-03}]{北科大資工所 113 學年度 第三題}
（是非；只寫答案）

Let $A_1$, \ldots, $A_6$ be matrices with dimension 5$\times$10, 10$\times$3, 3$\times$12, 12$\times$5,
5$\times$50, 50$\times$6, respectively. Let $C_{i,j}$ be the smallest number of scalar multiplications needed
for computing the matrix product $A_iA_{i+1}\ldots A_j$, assuming that multiplying an $r\times s$ matrix with
an $s\times t$ matrix takes $rst$ scalar multiplications. With the dynamic algorithm for Matrix-chain
Multiplication problem, mark by T(=true) or F(=false) each of the following statements: (8 pts)
\begin{enumerate}[label=(\arabic*)]
\item $C_{1,5}$ = 1655. (2 pts)
\item $C_{2,6}$ is computed before $C_{1,3}$. (2 pts)
\item $C_{2,5}$ is derived from $C_{2,3}$ and $C_{4,5}$. (2 pts)
\item The optimal order to multiply $A_1$, \ldots, $A_6$ with fewest number of scalar multiplications is
(($A_1A_2$)(($A_3$($A_4A_5$))$A_6$)) (2 pts)
\end{enumerate}
\end{pastexam}

\begin{pastsoln}
答案：(1) T　(2) F　(3) F　(4) F。

維度序列 $p = \langle 5, 10, 3, 12, 5, 50, 6\rangle$，$A_i$ 是 $p_{i-1} \times p_i$。 照資結書第~18.2~節：
\[
C_{i,i} = 0,\qquad C_{i,j} = \min_{i \le k < j}\bigl(C_{i,k} + C_{k+1,j} + p_{i-1}p_k p_j\bigr),
\]
依鏈長 $2, 3, \ldots, 6$ 的順序填表（資結書虛擬碼~18.4）。 每一格列出各個 $k$ 的值，取最小的（粗體）：
\begin{center}
\small
\begin{tabular}{lp{0.62\linewidth}l}
\toprule
格 & 各個 $k$ 的候選值 & $C_{i,j}$（最佳 $k$） \\
\midrule
長度 $2$ & $C_{1,2} = 150$，$C_{2,3} = 360$，$C_{3,4} = 180$，$C_{4,5} = 3000$，$C_{5,6} = 1500$ & \\
\midrule
$C_{1,3}$ & $k{=}1$: $0 + 360 + 5{\cdot}10{\cdot}12 = 960$；$k{=}2$: $150 + 0 + 5{\cdot}3{\cdot}12 = \mathbf{330}$ & $330$（2） \\
$C_{2,4}$ & $k{=}2$: $0 + 180 + 10{\cdot}3{\cdot}5 = \mathbf{330}$；$k{=}3$: $360 + 0 + 10{\cdot}12{\cdot}5 = 960$ & $330$（2） \\
$C_{3,5}$ & $k{=}3$: $0 + 3000 + 3{\cdot}12{\cdot}50 = 4800$；$k{=}4$: $180 + 0 + 3{\cdot}5{\cdot}50 = \mathbf{930}$ & $930$（4） \\
$C_{4,6}$ & $k{=}4$: $0 + 1500 + 12{\cdot}5{\cdot}6 = \mathbf{1860}$；$k{=}5$: $3000 + 0 + 12{\cdot}50{\cdot}6 = 6600$ & $1860$（4） \\
\midrule
$C_{1,4}$ & $580,\ \mathbf{405},\ 630$ & $405$（2） \\
$C_{2,5}$ & $\mathbf{2430},\ 9360,\ 2830$ & $2430$（2） \\
$C_{3,6}$ & $2076,\ \mathbf{1770},\ 1830$ & $1770$（4） \\
\midrule
$C_{1,5}$ & $4930,\ 1830,\ 6330,\ \mathbf{1655}$ & $1655$（4） \\
$C_{2,6}$ & $\mathbf{1950},\ 2940,\ 2130,\ 5430$ & $1950$（2） \\
\midrule
$C_{1,6}$ & $2250,\ \mathbf{2010},\ 2550,\ 2055,\ 3155$ & $2010$（2） \\
\bottomrule
\end{tabular}
\end{center}
（長度 $\ge 4$ 的格子，候選值依 $k = i, i+1, \ldots, j-1$ 列出。） 例如
$C_{1,5}$ 在 $k = 4$：$C_{1,4} + C_{5,5} + p_0 p_4 p_5 = 405 + 0 + 5 \cdot 5 \cdot 50 = 1655$；
$C_{2,5}$ 在 $k = 2$：$C_{2,2} + C_{3,5} + p_1 p_2 p_5 = 0 + 930 + 10 \cdot 3 \cdot 50 = 2430$。

\textbf{(1) T。} $C_{1,5} = 1655$，見上表。

\textbf{(2) F。} 由下往上的演算法依鏈長填表：$C_{1,3}$ 的鏈長是 $3$，$C_{2,6}$ 的鏈長是 $5$，$C_{2,6}$ 要用到 $C_{2,3}$、
$C_{2,4}$ 等較短的格子，一定在 $C_{1,3}$ \textbf{之後}才算。 （改用記憶化的由上往下遞迴，順序取決於呼叫的先後，
但題目說的是一般的動態規劃填表。）

\textbf{(3) F。} $C_{2,5}$ 的最佳切點是 $k = 2$，由 $C_{2,2}$ 和 $C_{3,5}$ 得到：$0 + 930 + 300 = 2430$。 用
$C_{2,3}$ 和 $C_{4,5}$（$k = 3$）得到的是 $360 + 3000 + 10 \cdot 12 \cdot 50 = 9360$，只是被比較過的候選值之一，最後的
值不是從它來的。

\textbf{(4) F。} 由切點表回推：$C_{1,6}$ 切在 $k = 2$，得 $(A_1A_2)(A_3 \cdots A_6)$；$C_{3,6}$ 切在 $k = 4$，得
$(A_3A_4)(A_5A_6)$。 最佳順序是
\[
\bigl((A_1A_2)\bigl((A_3A_4)(A_5A_6)\bigr)\bigr),\qquad 150 + 180 + 1500 + 3 \cdot 5 \cdot 6 + 5 \cdot 3 \cdot 6 = 2010.
\]
題目的 $((A_1A_2)((A_3(A_4A_5))A_6))$ 要 $A_4A_5$：$3000$；$A_3(A_4A_5)$：$3 \cdot 12 \cdot 50 = 1800$；再乘 $A_6$：
$3 \cdot 50 \cdot 6 = 900$；$A_1A_2$：$150$；最後 $5 \cdot 3 \cdot 6 = 90$；共 $5940$，比 $2010$ 多。
\end{pastsoln}
